\subsubsection{Simplex Tilings}

We connect simplex tiles together, via mutual reflection, to create a simplex tiling.
For the sake of brevity, when we say that a simplex tile reflects another tile (or that a tile is \emph{connected} to another tile), we mean that all simplex-chains in the first tile mutually reflect all simplex-chains in the second tile.
% There are an infinite number of simplex tilings, and an infinite number of \emph{patterns} by which we can construct simplex tilings.
A simplex tiling is thus a \emph{graph} of interconnected tiles.
Connected tiles are \emph{neighbors} and are \emph{adjacent} to one another in the tiling.

\defineTermTex{Simplex Tiling}{
    An interconnected graph of mutually reflecting simplex tiles
}

Before we look at some complex (and useful) tilings, let's consider three examples that are some of the simplest tilings possible: \autoref{fig:trivial-simplex-tiling}.
The examples contain 2, 3, and 4 simplex tiles respectively.

\aside{
    Note: we are not yet concerned with the security properties of these examples -- we need to build up to that first.
}

\begin{figure}
    \begin{subfigure}[m]{.32\textwidth}
        \vskip 0pt
        \centering
        \includegraphics[max height=\linewidth, max width=\linewidth]{tiling_alt_s5_d0_sag.pdf}
        \caption{2 simplex tiles}
        \label{fig:triv-tiling-2}
    \end{subfigure}%%
    \hfill
    \begin{subfigure}[m]{.32\textwidth}
        \vskip 0pt
        \centering
        \includegraphics[max height=\linewidth, max width=\linewidth]{tiling_a3_v3_s5_d0_sag.pdf}
        \caption{3 simplex tiles}
        \label{fig:triv-tiling-3}
    \end{subfigure}%%
    \hfill
    \begin{subfigure}[m]{.32\textwidth}
        \vskip 0pt
        \centering
        \includegraphics[max height=\linewidth, max width=\linewidth]{tiling_a4_v3_s5_d0_sag.pdf}
        \caption{4 simplex tiles}
        \label{fig:triv-tiling-4}
    \end{subfigure}
    \caption[Trivial simplex tilings]{Trivial simplex tilings -- these are all equivalent to a single simplex as each simplex-chain reflects all other simplex-chains in the network.}
    \label{fig:trivial-simplex-tiling}
\end{figure}

In those examples, notice that each simplex tile reflects each other tile.
This means that each simplex-chain is mutually reflecting \emph{all} other simplex-chains (in all other tiles), so each of these tilings is equivalent to a standalone simplex.
These tilings aren't particularly useful to us in practice, though, because we don't have a way to increase network-wide capacity beyond $O(c^2)$.

If we want $O(n)$ capacity, then we'll need to come up with a new pattern for how to connect tiles together and how to add new tiles.
Specifically: the pattern must \emph{always} have room for more simplex tiles, and each tile added must, on average, add \emph{more} that 1 new spot for additional tiles.\footnote{
    There are tilings where a tile consumes more than one free spot.
    For those tilings, on average each tile must add more spots than it consumes.
}

\subsubsection{Tile Valence}

The \emph{valence} ($v$) of a tile is the number of other tiles that it can connect to.
In our example earlier, we split our simplex tile's capacity for reflections into 4 equal parts -- each part was 12.5\% of each simplex-chain's \emph{total} capacity.
Since 1 of those parts is reserved for internal reflections, and 3 parts are reserved for external reflections, that simplex tile has a valence of $v=3$.

We don't \emph{have} to use a valence of 3 -- we can split up a tile's capacity for reflections any way that we like.
Each tile could even have different valencies, too, but this makes designing a tiling more complex.
For the sake of simplicity, we will only consider tilings where each tile has the same capacity and the same valence.
When we say that a \emph{tiling} (rather than a tile) has a valence of $v$, we mean that \emph{each tile in the tiling} has a valence of $v$.

If a tiling has a valence of $v=0$, then all capacity for reflections is reserved for internal reflections -- so this is just a normal simplex.

If a tiling has a valence of $v=1$, then the tiling is limited to 2 tiles at most -- there is only one solution.
We end up with a tiling that's equivalent to a normal simplex, though, since all simplex-chains \emph{must} reflect all other simplex-chains.
The complete tiling with $v=1$, which is the only possible tiling, is shown in \autoref{fig:triv-tiling-2}.

If a tiling has a valence of $v=2$, then there is a countably infinite number of solutions.
The trivial solution (which has 3 tiles) is where each simplex tile is connected to all other simplex tiles -- this is shown in \autoref{fig:triv-tiling-3}.
We can create a tiling of $v=2$ will 4 tiles, too: $A \lrar B \lrar C \lrar D \lrar A$.
This is part of a class of solutions: a loop.
With $v=2$, we can construct a loop of arbitrary length.
The last solution is a long chain of tiles that never loops back around; see \autoref{fig:tiling-v2-unbounded}.
This is our first \emph{unbounded} solution -- it's a solution where we can always add more tiles.
The problem with this solution is that, although it may have $O(n)$ capacity, the \emph{distance} between tiles is \emph{also} $O(n)$.
This is not practical since $O(n)$ SPV proofs would be needed to prove state on the far side of the tiling.
(The loop class of solutions is impractical for the same reason.)

\begin{figure}
    \centering
    \includegraphics[max width=.9\linewidth]{tiling_v2_s5_d6_sag.pdf}
    \caption[The only unbounded tiling at $v=2$.]{
        The only unbounded tiling at $v=2$.
        Although it is unbounded (we can always keep adding tiles), the \emph{distance} between those tiles is proportional to the size of the network -- so this isn't scalable.
    }
    \label{fig:tiling-v2-unbounded}
\end{figure}

Things start to get interesting when $v > 2$.

If a tiling has a valence of $v=3$, then there are many possible solutions (in fact, the number of solutions is uncountably infinite).
We will focus on one particular class: the patterns of tiling that are also \emph{trees}.

\subsubsection{Tree Tilings}

\begin{forest}
    [
        tree tiling
        [root tile]
        [first iteration]
        [further iterations]
    ]
\end{forest}

\begin{figure}
    \begin{subfigure}[m]{.32\textwidth}
        \vskip 0pt
        \centering
        \includegraphics[max width=\linewidth, max height=\linewidth]{tiling_s5_d1_sag}
        \caption{1st iteration. 4 tiles.}
        \label{fig:tree-tiling-v3-d1}
    \end{subfigure}%%
    \hfill
    \begin{subfigure}[m]{.32\textwidth}
        \vskip 0pt
        \centering
        \includegraphics[max width=\linewidth, max height=\linewidth]{tiling_s5_d2_sag}
        \caption{2nd iteration. 10 tiles.}
        \label{fig:tree-tiling-v3-d2}
    \end{subfigure}%%
    \hfill
    \begin{subfigure}[m]{.32\textwidth}
        \vskip 0pt
        \centering
        \includegraphics[max width=\linewidth, max height=\linewidth]{tiling_s5_d3_sag}
        \caption{3rd iteration. 22 tiles.}
        \label{fig:tree-tiling-v3-d3}
    \end{subfigure}
    \caption{}
    \label{fig:tree-tiling-v3}
\end{figure}

\subsubsubsection{\titlemath{$n$}{n}-ary Tree Tilings}

We will break our rules slightly: the root tile will have both twice the capacity for internal reflections than it would otherwise and one fewer children.
That is: the root tile reserves $2/v$ of it's capacity for reflections for internal reflections, instead of $1/v$ (as is the case for all other tiles), and it has $v-2$ children instead of $v-1$.
\autoref{fig:binary-tree-tiling} shows the first 3 iterations for $v=3$, and \autoref{fig:ternary-tree-tiling} shows the first 3 iterations for $v=4$.

\begin{figure}
    \begin{subfigure}[m]{.32\textwidth}
        \vskip 0pt
        \centering
        \includegraphics[max width=\linewidth, max height=\linewidth]{tiling_v3_nary_s5_d1_sag}
        \caption{1st iteration. 3 tiles.}
        \label{fig:binary-tree-tiling-d1}
    \end{subfigure}%%
    \hfill
    \begin{subfigure}[m]{.32\textwidth}
        \vskip 0pt
        \centering
        \includegraphics[max width=\linewidth, max height=\linewidth]{tiling_v3_nary_s5_d2_sag}
        \caption{2nd iteration. 7 tiles.}
        \label{fig:binary-tree-tiling-d2}
    \end{subfigure}%%
    \hfill
    \begin{subfigure}[m]{.32\textwidth}
        \vskip 0pt
        \centering
        \includegraphics[max width=\linewidth, max height=\linewidth]{tiling_v3_nary_s5_d3_sag}
        \caption{3rd iteration. 15 tiles.}
        \label{fig:binary-tree-tiling-d3}
    \end{subfigure}
    \caption{}
    \label{fig:binary-tree-tiling}
\end{figure}

\begin{figure}
    \begin{subfigure}[m]{.32\textwidth}
        \vskip 0pt
        \centering
        \includegraphics[max width=\linewidth, max height=\linewidth]{tiling_v4_nary_s5_d1_sag}
        \caption{1st iteration. 4 tiles.}
        \label{fig:ternary-tree-tiling-d1}
    \end{subfigure}%%
    \hfill
    \begin{subfigure}[m]{.32\textwidth}
        \vskip 0pt
        \centering
        \includegraphics[max width=\linewidth, max height=\linewidth]{tiling_v4_nary_s5_d2_sag}
        \caption{2nd iteration. 13 tiles.}
        \label{fig:ternary-tree-tiling-d2}
    \end{subfigure}%%
    \hfill
    \begin{subfigure}[m]{.32\textwidth}
        \vskip 0pt
        \centering
        \includegraphics[max width=\linewidth, max height=\linewidth]{tiling_v4_nary_s5_d3_sag}
        \caption{3rd iteration. 40 tiles.}
        \label{fig:ternary-tree-tiling-d3}
    \end{subfigure}
    \caption{}
    \label{fig:ternary-tree-tiling}
\end{figure}

\subsubsection{Tree Tilings: Preliminary Scaling Analysis}

\newpage
